% RESEARCH_DRAFT_SHA256: ca018d944639258ecc811afe9c10c5f3293aa8808ca97291ab0a40cbdbf2c4bb
% Layout and source-reference wrappers expanded for the web renderer.
\documentclass[11pt]{article}
\usepackage[margin=25mm]{geometry}
\usepackage{amsthm}
\usepackage{fontspec,unicode-math}
\setmainfont{Latin Modern Roman}
\setmathfont{Latin Modern Math}
\usepackage{amsmath,amssymb,mathtools,xurl,hyperref}
\hypersetup{colorlinks=true,urlcolor=blue,linkcolor=blue}
\setcounter{secnumdepth}{0}
\setlength{\parindent}{0pt}
\setlength{\parskip}{5pt}
\setlength{\emergencystretch}{3em}
\newcommand{\sref}[2]{\hyperlink{source:#1:#2}{[S#2]}}
\newcommand{\cataloguescope}[1]{\paragraph{Source record.}#1}
\newcommand{\floor}[1]{\left\lfloor#1\right\rfloor}
\newtheorem{proposition}{Proposition}
\newtheorem{lemma}[proposition]{Lemma}
\begin{document}
\paragraph{Research manuscript.} OpenAI \texttt{gpt-6-astra}, reasoning \texttt{ultra}; 11 October 2026.
\paragraph{Preserved source body.} The complete catalogue mathematical text and sources follow with layout and source links expanded, including source qualifications and its historical preparation notes. New research begins at the section headed Research attempt.

\section{A co-governed category-and-access charter beyond a single gender binary}\label{statement}

\subsection{Definitions and mathematical statement}

\paragraph{Shared notation.}Standard number systems are $\mathbb N,\mathbb Z,\mathbb Q,\mathbb R,\mathbb C$. Specialized notation and index ranges follow the cited source and the definitions below.


Gender diversity describes participation across differing gender identities. Organizational membership criteria determine who a named organization includes, while inclusive practice concerns how people can participate. A binary classification has two categories and may fail to represent the categories relevant to the intended community.
\paragraph{Record scope and qualifications.}
The catalogue title is retained as metadata. Named certificates, contracts, bounds or frameworks in that title are not asserted as results of the original source extract.
No particular organization, membership rule, or measured outcome is identified. The extract is a community-policy discussion rather than a mathematical assertion.

Catalogue screening limitation: Community, history, equity, or organizational policy discussion, rather than a mathematical statement.

The source snapshot labels this record partially solved; this is catalogue metadata rather than independent verification of a mathematical result or the current completion of a project.
\paragraph{Statement recorded in the supplied source.}
Gender Diversity Beyond the Gender Binary

How can we address and work towards gender diversity outside of the gender binary? Can "Women in Math" organizations clarify what they mean by "women?" And can we push for greater inclusivity in these organizations? \href{http://aimpl.org/gemscombin/7/}{[S2]}
\paragraph{Source status.}
The supplied catalogue labels this source record partially solved; that metadata was checked against the supplied snapshot on 2026-10-11. The label is preserved for provenance and is not an independent verification of present open status or mathematical correctness.
\subsection{Short English statement}
Consider gender diversity beyond binary classifications and clarify inclusive participation in mathematics organizations.
\subsection{Sources}
\begin{itemize}
\item[S1] Ulam AI and the UnsolvedMath Contributors, UnsolvedMath: A Curated Collection of Open Mathematics Problems, catalogue entry. Catalogue curation and metadata: CC BY 4.0. The source extract is retained; definitions, English summaries and qualifications are editorial additions.
\url{https://huggingface.co/datasets/ulamai/UnsolvedMath}.
\item[S2] AIM Workshop Problem Lists, source linked by the catalogue.
\url{http://aimpl.org/gemscombin/7/}.
\end{itemize}
\paragraph{Preparation and review.}
The original source and UnsolvedMath contributors retain attribution for the source extract. Codex (GPT-6; exact serving revision not exposed) assisted with definitions, English summaries, scope notes and formatting. No new human mathematical authorship or independently verified solution is asserted. Catalogue identity checked against the complete supplied registry and current repository. Distinct upstream records, including repeated questions and broad or damaged extracts, are retained for complete coverage of the source CS category. A repeated record is not a new independent conjecture. Complete mathematical-equivalence and current-literature review has not been performed. Preparation and final review follow contribution prompts C and D, with source gaps disclosed. Structural and rendering checks do not establish a new theorem or open status.
\label{end:statement}

\section{Research attempt}
\subsection{Literature, source context and scope}
The complete catalogue body above retains the original agenda, source status, and historical preparation. The AIMPL problem page could not be retrieved. The supplied complete source transcription was read, and the accessible primary AIM workshop announcement was checked for context. The announcement is not a substitute for verification of the inaccessible problem page. The catalogue's earlier proposed frameworks and research summaries are context, not established solutions credited to this attempt. The following is an explicitly limited investigation of one route; it neither substitutes a formalized toy model for the entire source agenda nor claims empirical or deployed-software results. Elementary auxiliary facts are not new original-target progress.

\subsection{Testing whether marginal category counts determine overlap}
The source calls for greater inclusivity and clarification of organizational categories. We examine only a data-representation pitfall: replacing voluntarily overlapping membership categories by marginal totals and then inferring their intersection. Let $A,B$ be two possibly overlapping self-described categories in a finite universe of size $N$, with $|A|=a$, $|B|=b$. The names of categories and who elects to use them are not assigned by this model.
\begin{proposition}
The possible intersection sizes $r=|A\cap B|$ are exactly the integers satisfying
\[
\max(0,a+b-N)\le r\le\min(a,b).
\]
\end{proposition}
\begin{proof}
The intersection is a subset of each set and has nonnegative size, yielding the upper bound and zero lower bound. Since $|A\cup B|=a+b-r\le N$, also $r\ge a+b-N$. Conversely, for any integer in the stated interval, form four disjoint cells of sizes $r,a-r,b-r,N-a-b+r$. The inequalities say precisely that all four are nonnegative, and their sum is $N$. Define $A$ as the first two cells and $B$ as the first and third. They have the specified marginals and intersection, proving attainability.
\end{proof}
With $N=100,a=60,b=50$, any intersection size from 10 through 50 is compatible with the same marginals. A system that makes the categories mutually exclusive would silently rule out every one of these possible configurations. This demonstrates a mathematical information loss in a specified schema; it does not decide how any real organization should define membership or how individuals should identify.
\subsection{Verification and remaining gaps}
All 41 feasible intersections were checked by constructing their four cell sizes. No organizational charter, consultation, category definition or inclusion outcome was delivered. The sharp set-theoretic bounds are standard auxiliary facts, not a solution to the community agenda.

\section{References and source attribution}
Primary source: \url{http://aimpl.org/gemscombin/7/}. Attempted primary AIMPL page retrieval; failed with timeout or gateway error. The supplied full source transcription was read; the inaccessible live page was not independently inspected.

\url{https://aimath.org/pastworkshops/gemscombin.html}. Read the primary March 2023 workshop announcement and its community inclusion agenda.

All retrievals were checked on 11 October 2026. Original contributor and catalogue attribution remains in the preserved Sources section. The elementary arguments below are self-contained; their exposition is not a claim of first discovery.


\paragraph{Integration review.} Codex (GPT-6; exact serving revision not exposed) reviewed the written arguments and expanded layout and source-reference wrappers. Repository metadata and the production web renderer were checked using Python 3.13 and JavaScript with MathJax 3.2.2. This model-assisted review is not independent human proof verification; no PDF compilation or formal proof checking is asserted.

\section{Completion Estimate}
\noindent\textbf{COMPLETION: 0\%}
No substantiated new progress toward the entire original target is claimed. The worked results and failed route are delimited above; a known result, illustrative model, or auxiliary obstruction is not counted as original-target completion.

\section{Final status and attribution}
\textbf{Attempt status: unresolved.} This describes this research attempt, not a new ruling on the source catalogue's status or on external literature claims.
The mathematics and this self-audit were generated by OpenAI \texttt{gpt-6-astra}, reasoning \texttt{ultra}, on 11 October 2026. Codex (GPT-6; exact serving revision not exposed) orchestrated the batch and contributed formatting/repository review. No human mathematical authorship, independent human proof checking, or formal verification is claimed. The historical preparation paragraphs in the reproduced catalogue body concern that body only.

Contribution prompt A was applied to the complete definition and supplied source record, followed by the adversarial checks of prompt B. Local reproduction files are \texttt{batch013/checks.py}, \texttt{computations.json}, and the per-ID audit JSON. They were not public links at preparation time. Numerical checks are finite and do not prove unrestricted results. No LaTeX engine was found locally; compilation was not performed by this research agent.

\end{document}
